% exact-proofs.tex -- exact output, finite recovery without uniqueness,
% explicit polynomial factors. Fragment: no preamble.
% Uses the shared notation of the paper. Cross-references to other sections:
%   def:curvature (L_i), eq:growth / def:growth (point growth, kappa),
%   eq:step (graded step rule), FG / CT and thm:approx (certified approximation),
%   rem:termination (CT stops on every instance), lem:dp (two-pass DP).

\section{Exact output}\label{sec:exact}

A certified gap below $2^{-q}$ is not an exact solution. For rational
quadratics, a global minimizer and the optimal value have bounded rational
height, and this makes a sufficiently small certified gap convertible into an
exact optimizer that is checked against the original data. This section
proves the height bound in a form that serves two purposes at once, gives an
acceptance rule that needs neither uniqueness nor growth, and shows that an
exact optimizer can be recovered from any feasible point that is close to the
\emph{set} of minimizers. Under quadratic growth at a unique minimizer this
yields exact output within the bound of Theorem~\ref{thm:approx}. Without
uniqueness, the single-center grids of Section~\ref{sec:growth} can be
exponentially slow, and we show that unions of uniform cells restore a rate
that is polynomial for fixed bag size. The last subsection treats explicit
polynomial factors.

\subsection{Scaled data}\label{sec:exact-data}

Throughout Sections~\ref{sec:exact-data}--\ref{sec:exact-nonunique}, $F$ is a
quadratic with rational data,
\[
F(x)=\tfrac12x^{\top}Hx+b^{\top}x+c ,
\]
where $H$ is the symmetric Hessian, on a bounded mixed box
$X=\prod_iX_i$ with rational endpoints, after fixed coordinates have been
substituted and integer endpoints rounded inward. Write
$S=\argmin_XF$ and $\OPT=\min_XF$. Let $D$ be a positive common denominator
of the monomial coefficients of $F$, namely $H_{ii}/2$, $H_{ij}$ ($i<j$),
$b_i$ and $c$, and of all endpoints $\ell_i,u_i$. Put
\begin{equation}\label{eq:exact-constants}
P=DH,\qquad I_C^+=\{i\in I_C:H_{ii}>0\},\qquad
R=D\prod_{i\in I_C^+}P_{ii},\qquad V=DR^2,\qquad \tau=\frac1{4nR}.
\end{equation}
The matrix $P$ is integral and symmetric, since $P_{ii}=2D(H_{ii}/2)$ and
$P_{ij}=DH_{ij}$; this choice of $D$ absorbs the factor two of the
off-diagonal entries automatically. Each $P_{ii}$ with $i\in I_C^+$ is a
positive integer, and $P_{ii}\le DL$ with $L=\max_iL_i$, so
$\log_2R\le\log_2D+\abs{I_C^+}\log_2(DL)$ is polynomial in $I$. The numbers
$R$, $V$ and $\tau$ are computed from the original data; refined boxes never
enter them.

Two elementary facts are used repeatedly. If every coordinate of $x\in X$
lies in $\rho^{-1}\Z$ for a positive integer $\rho$, then
$D\rho^2F(x)\in\Z$, because $D$ clears every monomial coefficient. If
$a/W\ne a'/W'$ are rationals with positive denominators, then
$\abs{a/W-a'/W'}\ge1/(WW')$.

\begin{lemma}[Hadamard's inequality]\label{lem:hadamard}
If $A$ is a real symmetric positive definite $m\times m$ matrix, then
$\det A\le\prod_{i=1}^mA_{ii}$.
\end{lemma}

\begin{proof}
Let $\Lambda=\diag(A_{11},\dots,A_{mm})$, which is positive definite, and
$B=\Lambda^{-1/2}A\Lambda^{-1/2}$. Then $B$ is positive definite with unit
diagonal, so its eigenvalues $\lambda_k>0$ satisfy $\sum_k\lambda_k=m$, and
by the arithmetic--geometric mean inequality
$\det B=\prod_k\lambda_k\le1$. Hence
$\det A=\det B\prod_iA_{ii}\le\prod_iA_{ii}$.
\end{proof}

\subsection{Stationary polytopes and rational height}

\begin{definition}[Stationary polytope]\label{def:statpoly}
For $s\in S$ let $J_0(s)=\{i\in I_C:\ell_i<s_i<u_i\}$ and
\[
\mathcal P(s)=\bigl\{x\in\R^n:\ x_i=s_i\ (i\notin J_0(s)),\ \
\ell_i\le x_i\le u_i\ \text{and}\ \partial_iF(x)=0\ (i\in J_0(s))\bigr\}.
\]
\end{definition}

\begin{lemma}[Stationary polytopes]\label{lem:statpoly}
Let $s\in S$ and $J_0=J_0(s)$.
\begin{enumerate}[label=(\alph*)]
\item $H_{J_0J_0}\succeq0$.
\item $\mathcal P(s)$ is a nonempty polytope contained in $X$, and
$\mathcal P(s)\subseteq S$.
\item Let $v$ be a vertex of $\mathcal P(s)$ and
$T=\{i\in J_0:\ell_i<v_i<u_i\}$. Then $P_{TT}\succ0$, every coordinate of $v$
lies in $(D\det P_{TT})^{-1}\Z$, and $1\le D\det P_{TT}\le R$ (with
$\det P_{TT}=1$ if $T=\emptyset$).
\end{enumerate}
\end{lemma}

\begin{proof}
(a) For $w\in\R^n$ supported on $J_0$, the point $s+tw$ lies in $X$ for all
small $\abs t$, because the coordinates in $J_0$ are continuous and strictly
inside their bounds. Since $F(s+tw)=\OPT+t\nabla F(s)^{\top}w+\frac{t^2}2w^{\top}Hw$
has a minimum at $t=0$, we get $\nabla_{J_0}F(s)=0$ and $w^{\top}Hw\ge0$.

(b) By (a), $s\in\mathcal P(s)$. The set is defined by finitely many linear
equations and inequalities inside a bounded box, so it is a polytope. Its
points keep the integer coordinates of $s$ and satisfy all bounds, so
$\mathcal P(s)\subseteq X$. For $x\in\mathcal P(s)$, $d=x-s$ is supported on
$J_0$, and $H_{J_0J_0}d_{J_0}=\nabla_{J_0}F(x)-\nabla_{J_0}F(s)=0$. Hence
\[
F(x)-F(s)=\nabla F(s)^{\top}d+\tfrac12d^{\top}Hd
=\nabla_{J_0}F(s)^{\top}d_{J_0}+\tfrac12d_{J_0}^{\top}H_{J_0J_0}d_{J_0}=0 ,
\]
so $x\in S$.

(c) Suppose $w\in\R^T$ is nonzero and $H_{J_0T}w=0$. Extend $w$ by zero to
$\tilde w\in\R^n$. Then $v\pm t\tilde w\in\mathcal P(s)$ for small $t>0$: the
equations are preserved, the coordinates in $T$ have slack, and all other
coordinates are unchanged. This contradicts that $v$ is a vertex. So
$H_{J_0T}$ has trivial kernel. If $w\in\R^T\setminus\{0\}$ had
$w^{\top}H_{TT}w=\tilde w^{\top}H\tilde w=0$, positive semidefiniteness of
$H_{J_0J_0}$ would give $H_{J_0J_0}\tilde w_{J_0}=0$, that is
$H_{J_0T}w=0$; hence $H_{TT}\succ0$ and $P_{TT}=DH_{TT}\succ0$. In
particular $H_{ii}>0$ for $i\in T$, so $T\subseteq I_C^+$.

The coordinates of $v$ outside $T$ are integers or endpoints, so they lie in
$D^{-1}\Z$. The rows $i\in T$ of $\nabla F(v)=0$, multiplied by $D^2$, read
\[
P_{TT}(Dv_T)=-D(Db_T)-\sum_{j\notin T}P_{Tj}(Dv_j),
\]
whose right-hand side is integral. Cramer's rule gives
$Dv_T\in(\det P_{TT})^{-1}\Z^T$, so every coordinate of $v$ lies in
$(D\det P_{TT})^{-1}\Z$. Finally $\det P_{TT}$ is a positive integer, and by
Lemma~\ref{lem:hadamard} and $T\subseteq I_C^+$,
$\det P_{TT}\le\prod_{i\in T}P_{ii}\le\prod_{i\in I_C^+}P_{ii}$.
\end{proof}

\begin{corollary}[Rational height]\label{cor:height}
(a) Some global minimizer has all coordinates in $\rho^{-1}\Z$ for a positive
integer $\rho\le R$; every isolated global minimizer has this property.
(b) $\OPT=a/W$ in lowest terms with $1\le W\le V$.
\end{corollary}

\begin{proof}
(a) Take any $s\in S$. The nonempty polytope $\mathcal P(s)$ has a vertex,
which is a global minimizer by Lemma~\ref{lem:statpoly}(b) and has the
stated height by Lemma~\ref{lem:statpoly}(c). If $s$ is isolated in $S$, then
the convex set $\mathcal P(s)\subseteq S$ contains no other point near $s$,
hence equals $\{s\}$, and $s$ is its vertex.
(b) Evaluate $F$ at the minimizer of (a): $D\rho^2\OPT\in\Z$, so
$W\le D\rho^2\le V$.
\end{proof}

\begin{remark}[Other admissible height constants]\label{rem:heights}
Any bound on the minors of $P$ restricted to $I_C$ also works in
Lemma~\ref{lem:statpoly}(c) and Corollary~\ref{cor:height}, for example
$\abs{I_C}!\,(\max_{ij}\abs{P_{ij}})^{\abs{I_C}}$ by the Leibniz expansion, or
$\prod_{i\in I_C}\max\{1,\sum_{j\in I_C}\abs{P_{ij}}\}$ by Hadamard's
inequality for rows. Both also bound \emph{nonprincipal} minors, which arise
if a vertex of $\mathcal P(s)$ is described by a nonprincipal set of
stationarity rows. Part~(c) shows that a principal system $P_{TT}$ always
suffices, and Lemma~\ref{lem:hadamard} then gives the diagonal bound
in~\eqref{eq:exact-constants}, which is never larger than either
alternative and depends only on the coordinate curvatures. All these
constants have polynomial bit length; only the threshold values below
change.
\end{remark}

\subsection{Acceptance and recovery}

\begin{proposition}[Exact acceptance]\label{prop:accept}
Let $\beta\le\OPT$ and let $x\in X$ with $F(x)=a/W$ in lowest terms. If
$F(x)-\beta<1/(VW)$, then $x\in S$.
\end{proposition}

\begin{proof}
If $F(x)>\OPT$, then $F(x)-\OPT\ge1/(VW)$, since $\OPT$ has denominator at
most $V$ by Corollary~\ref{cor:height}(b). With $\beta\le\OPT$ this
contradicts $F(x)-\beta<1/(VW)$.
\end{proof}

The rule uses the denominator of the candidate itself. It does not assume
that feasible values lie on a lattice, which is false when continuous
coordinates vary.

\paragraph{Procedure REC$(y)$ (snapping recovery).}
Given $y\in X$: fix every integer coordinate at $y_i$; for every continuous
coordinate, fix it at $\ell_i$ if $y_i-\ell_i\le\tau$, else at $u_i$ if
$u_i-y_i\le\tau$, and otherwise call it free; let $J$ be the free set. Find
by exact linear programming a point $\hat x$ that agrees with the fixed
values, satisfies $\ell_J\le\hat x_J\le u_J$ and $\nabla_JF(\hat x)=0$, or
report failure. The data of this linear system are original coefficients,
endpoints and integer values of $y$, so it is solved in time polynomial in
$I$ \cite[Theorem~6.4.12]{GrotschelLovaszSchrijver1988}. If $H_{JJ}$ is
nonsingular, Gaussian elimination suffices.

\begin{lemma}[Snapping recovery]\label{lem:snap}
If $y\in X$ and $\dist(y,S)\le\tau/2$, then REC$(y)$ does not fail, and every
point it can return lies in $S$.
\end{lemma}

\begin{proof}
Let $s\in S$ be nearest to $y$; it exists since $S$ is compact. As
$\norm{y-s}\le\tau/2<1$, the integer coordinates of $y$ and $s$ agree. Two
distinct endpoints of a continuous coordinate differ by at least
$1/D\ge4n\tau>2\tau$, so no coordinate is within $\tau$ of both endpoints. If
$s_i=\ell_i$ then $y_i-\ell_i\le\tau/2$ and coordinate $i$ is fixed at
$\ell_i$; similarly for $s_i=u_i$. Hence every continuous coordinate at a
bound in $s$ is fixed at its value in $s$, and $J\subseteq J_0=J_0(s)$. Let
$E=J_0\setminus J$ be the remaining fixed coordinates, each fixed at an
endpoint $e_i\in\{\ell_i,u_i\}$, and define on $\mathcal P(s)$ the linear
function
\[
\phi(x)=\sum_{i\in E}\abs{x_i-e_i}
=\sum_{i\in E,\,e_i=\ell_i}(x_i-\ell_i)+\sum_{i\in E,\,e_i=u_i}(u_i-x_i)\ge0 .
\]
For $i\in E$ we have $\abs{y_i-e_i}\le\tau$ and $\abs{s_i-y_i}\le\tau/2$, so
$\phi(s)\le\frac32\abs E\tau\le\frac3{8R}<\frac1R$. The minimum of $\phi$
over the polytope $\mathcal P(s)$ is attained at a vertex $v$. By
Lemma~\ref{lem:statpoly}(c) the coordinates of $v$ and all endpoints lie in
$\rho^{-1}\Z$ for some integer $\rho\le R$, so $\phi(v)\in\rho^{-1}\Z$.
Since $0\le\phi(v)\le\phi(s)<1/R\le1/\rho$, we get $\phi(v)=0$.

Thus $s'=v\in\mathcal P(s)\subseteq S$ satisfies $s'_i=e_i$ for $i\in E$,
agrees with $s$ (hence with the fixed values of REC) outside $J_0$, and has
$\nabla_JF(s')=0$ because $J\subseteq J_0$. So the system solved by REC is
feasible. If $\hat x$ is any solution, then $d=\hat x-s'$ is supported on
$J$ and $H_{JJ}d_J=\nabla_JF(\hat x)-\nabla_JF(s')=0$, so
$F(\hat x)-F(s')=\nabla_JF(s')^{\top}d_J+\frac12d_J^{\top}H_{JJ}d_J=0$.
Since $\hat x\in X$, it is a global minimizer. No nonsingularity of $H_{JJ}$
is used.
\end{proof}

\begin{theorem}[From certified gaps to exact minimizers]\label{thm:transfer}
Let $\beta\le\OPT$ be certified and $y\in X$.
\begin{enumerate}[label=(\alph*)]
\item If $F$ has set growth $F(x)-\OPT\ge g_S\dist(x,S)^2$ on $X$ with
$g_S>0$, and
\begin{equation}\label{eq:transfer-threshold}
F(y)-\beta\le\varepsilon_S:=\min\Bigl\{\frac{g_S\tau^2}{4},\
\frac1{2V^2}\Bigr\}
=\min\Bigl\{\frac{g_S}{64n^2R^2},\ \frac1{2V^2}\Bigr\},
\end{equation}
then REC$(y)$ returns a point $\hat x$ that passes the test of
Proposition~\ref{prop:accept}.
\item Without any growth assumption there is $\varepsilon_0>0$, depending on
the instance, such that the conclusion of (a) holds whenever
$F(y)-\beta\le\varepsilon_0$.
\end{enumerate}
\end{theorem}

\begin{proof}
(a) $g_S\dist(y,S)^2\le F(y)-\OPT\le F(y)-\beta\le g_S\tau^2/4$, so
Lemma~\ref{lem:snap} applies and $\hat x\in S$. Then
$F(\hat x)=\OPT=a/W$ with $W\le V$, and
$F(\hat x)-\beta\le F(y)-\beta\le1/(2V^2)<1/(VW)$.
(b) The set $Y=\{x\in X:\dist(x,S)\ge\tau/2\}$ is compact and $F>\OPT$ on
it, so $\eta=\min_YF-\OPT>0$ (put $\eta=\infty$ if $Y=\emptyset$). With
$\varepsilon_0=\min\{\eta/2,1/(2V^2)\}$, $F(y)-\beta\le\varepsilon_0$ gives
$\dist(y,S)<\tau/2$, and the argument of (a) applies.
\end{proof}

\begin{remark}[Set growth is automatic]\label{rem:setgrowth}
For every rational mixed box QP some $g_S>0$ exists. On an integer slice
whose continuous box contains a minimizer, Luo and Sturm's error bound for a
quadratic on a polytope \cite[Theorem~3.3]{LuoSturm2000}, applied to $F-\OPT$,
gives $F(x)-\OPT\ge c\,\dist(x,S\cap\text{slice})^2\ge c\,\dist(x,S)^2$; on
the finitely many other slices $F-\OPT$ is bounded below by a positive
constant while $\dist(x,S)$ is bounded by the diameter of $\bar X$. As for
point growth, no useful lower bound on $g_S$ follows. Theorem~\ref{thm:transfer}
shows that exactness costs only $\log_2(1/\varepsilon_S)
=\poly(I)+\log_2(1/g_S)$ bits of certified accuracy; the cost of reaching that
accuracy is a property of the approximation algorithm.
\end{remark}

\subsection{Exact output under quadratic growth}

\begin{lemma}[Uniqueness implies growth]\label{lem:unique-growth}
If $S=\{x^*\}$, then~\eqref{eq:growth} holds for some $g>0$.
\end{lemma}

\begin{proof}
Otherwise there are $x^k\in X\setminus\{x^*\}$ with
$(F(x^k)-\OPT)/\norm{x^k-x^*}^2\to0$. As $\bar X$ is bounded,
$F(x^k)\to\OPT$, and compactness of $X$ with uniqueness gives $x^k\to x^*$;
so the integer coordinates of $x^k$ equal those of $x^*$ for large $k$;
since $x^k\ne x^*$, they differ in some continuous coordinate. Put
$t_k=\norm{x^k-x^*}$ and
$u^k=(x^k-x^*)/t_k$; pass to a subsequence with $u^k\to u$, $\norm u=1$,
$u_{I_Z}=0$. Then
\[
\frac{F(x^k)-\OPT}{t_k^2}=\frac{\nabla F(x^*)^{\top}u^k}{t_k}
+\frac12(u^k)^{\top}Hu^k .
\]
Since $x^*$ minimizes $F$ over the convex continuous box of its integer slice
and $x^k$ lies in that box, the first term is nonnegative. The left side
tends to zero and the second term is bounded, so
$\nabla F(x^*)^{\top}u=0$ and $u^{\top}Hu\le0$. Each $u^k$ satisfies
$u^k_i\ge0$ when $x^*_i=\ell_i$ and $u^k_i\le0$ when $x^*_i=u_i$; so does
$u$, and hence $x^*+tu\in X$ for small $t>0$. Then
$F(x^*+tu)-\OPT=\frac{t^2}2u^{\top}Hu\le0$, contradicting uniqueness.
\end{proof}

The constant $g$ can be exponentially small in $I$.

\paragraph{Algorithm EX (exact output).}
For $q=1,2,4,8,\dots$: run CT$(2^{-q})$ and obtain a feasible $\hat x$ and a
certified $\beta$; run REC$(\hat x)$; if it returns $\tilde x$ with
$F(\tilde x)-\beta<1/(VW)$, where $W$ is the reduced denominator of
$F(\tilde x)$, return $\tilde x$, $F(\tilde x)$ and the certificate of
CT. If $L=0$, CT is exact at its first stage and the test passes with
$\tilde x=\hat x$.

\begin{theorem}[Exact rational output]\label{thm:exact}
For every rational mixed box QP with a supplied tree decomposition, algorithm
EX stops, and its output is a global minimizer and $\OPT$, with a certificate
whose validity uses neither growth nor uniqueness. If the instance has
quadratic growth~\eqref{eq:growth} with condition number $\kappa$, EX uses
at most $f_1(p,\kappa)(I+1)^{C_1}$ bit operations, where $C_1$ is an absolute
constant; neither $g$ nor $\kappa$ is an input. This bound remains valid if
REC skips
every call whose matrix $H_{JJ}$ is singular, so under growth no linear
programming is needed.
\end{theorem}

\begin{proof}
\emph{Validity.} Each returned point passes Proposition~\ref{prop:accept}
with a certified $\beta$ (Theorem~\ref{thm:approx}), so it is a global
minimizer.

\emph{Termination.} CT stops on every instance
(Remark~\ref{rem:termination}). Once $2^{-q}\le\varepsilon_0$ of
Theorem~\ref{thm:transfer}(b), the test passes.

\emph{Bound under growth.} Point growth is set growth with $S=\{x^*\}$ and
$g_S=g$. By Theorem~\ref{thm:transfer}(a), EX stops at the first $q$ with
$2^{-q}\le\varepsilon_S$, so the largest $q$ used is less than
$2q^*$ with $q^*=\lceil\log_2(1/\varepsilon_S)\rceil$. Since $g\ge L/\kappa$
and $\log R$, $\log V$, $\log n$, $\log(1/L)$ are at most polynomial in $I$,
$q^*\le\poly(I)+\log_2\kappa$. There are $O(\log q^*)$ calls, each within
$f(p,\kappa)(I+2q^*+1)^C$ bit operations by Theorem~\ref{thm:approx}, and
$(I+2q^*+1)^C\le(1+\log_2\kappa)^{C}\poly(I)$. REC, the evaluation of
$F(\tilde x)$ and the comparison cost $\poly(I+q^*)$. Absorbing all powers
of $1+\log\kappa$ into $f_1$ proves the bound. Finally, if
$\dist(\hat x,x^*)\le\tau/2$, then $J\subseteq J_0(x^*)$ as in
Lemma~\ref{lem:snap}, and $H_{J_0J_0}\succ0$: a nonzero $w$ with
$H_{J_0J_0}w=0$ would make $F$ constant on a short segment through $x^*$.
So $H_{JJ}$ is nonsingular at every call that the argument above needs, and
skipping singular calls does not change the bound.
\end{proof}

\begin{remark}[Coordinatewise reconstruction]\label{rem:cf}
Under uniqueness one may instead round each coordinate of $\hat x$ to the
unique rational of denominator at most $R$ within $1/(4R^2)$ (continued
fractions, \cite[Theorem~5.1.9]{GrotschelLovaszSchrijver1988}); this succeeds
once $2^{-q}\le g/(32R^4)$, by Corollary~\ref{cor:height}(a) and the
$1/R^2$ separation of such rationals. Snapping recovery needs only
$g/(64n^2R^2)$, does not need uniqueness, and returns a minimizer even when
$S$ is a continuum.
\end{remark}

\subsection{Nonunique minimizers}\label{sec:exact-nonunique}

Theorem~\ref{thm:transfer} reduces exact output to reaching accuracy
$\varepsilon_S$. The algorithm of Section~\ref{sec:growth} does not reach it
at a rate controlled by $(p,\kappa)$ when the minimizer is not unique, even
with two isolated minimizers and set growth: its grids are graded around one
center, and coarse cells near a distant minimizer keep the corrected minimum
low.

\begin{proposition}[One center cannot certify two distant minimizers]
\label{prop:twocenters}
For $M>0$ let
\[
F_M(x,z)=x^2-2xz+Mz\qquad\text{on }X=[0,M]^2\ (\text{both coordinates
continuous}).
\]
Then $S=\{(0,0),(M,M)\}$, $\OPT=0$, $L_x=2$, $L_z=0$, and
$F_M(v)\ge\frac1{20}\dist(v,S)^2$ on $X$, so the set-growth condition number
is at most $40$; one bag $\{x,z\}$ covers both scopes ($p=2$).
Let $\theta\in(0,1/4]$, $h\in(0,M]$ and let $c\in X$ be any center. For the
graded grids~\eqref{eq:step} on $X$ with center $c$, mesh $h$, grading
$\theta$, and any corrections $d_x=L_xw_x^2/8$ with $L_x\ge2$ and $d_z\ge0$,
\begin{equation}\label{eq:twocenters}
\beta=\min_GQ\le-\max\Bigl\{\frac{\theta^2M^2}{379},\ \frac{h^2}{20}\Bigr\}.
\end{equation}
Consequently, in FG and CT every stage has certified gap
$U-\beta\ge\max\{\theta^2M^2/379,\,h^2/20\}$; a stage certifying a
gap $\varepsilon$ uses grids with at least $M/(48\sqrt\varepsilon)$ nodes in
each coordinate; and any exact-output procedure based on these certificates
that accepts only when $U-\beta<1$, such as EX, builds a table with at least
$M^2/2304$ entries. For integer $M$ the input length is $I=O(\log M)$, so
this work is $2^{\Omega(I)}$.
\end{proposition}

\begin{proof}
\emph{The instance.} $F_M=x^2+z(M-2x)$. For $x<M/2$ the minimum over $z$ is
at $z=0$ with value $x^2$; for $x>M/2$ it is at $z=M$ with value $(M-x)^2$;
at $x=M/2$ the value is $M^2/4$. So $S$ and $\OPT$ are as stated, and the
Hessian $\bigl(\begin{smallmatrix}2&-2\\-2&0\end{smallmatrix}\bigr)$ gives
$L_x=2$, $L_z=0$. The map $(x,z)\mapsto(M-x,M-z)$ preserves $F_M$ (a direct
expansion) and swaps the minimizers, so for growth we may assume
$x\le M/2$. If $x\le M/4$, then $M-2x\ge M/2\ge z/2$ and
$F_M\ge x^2+z^2/2\ge\frac12\norm{(x,z)}^2$. If $M/4<x\le M/2$, then
$F_M\ge x^2\ge M^2/16$ while $\norm{(x,z)}^2\le5M^2/4$. In both cases
$F_M\ge\frac1{20}\dist((x,z),S)^2$.

\emph{The bound~\eqref{eq:twocenters}.} By symmetry assume $M-c_x\ge M/2$
(otherwise use the reflected argument with $z=0$ and the endpoint $0$). Let
$c_x=v_0<v_1<\dots<v_K=M$ be the $x$-nodes on the right of the center,
$D_k=v_k-c_x$, $a=v_{K-1}-v_{K-2}$ (if $K\ge2$) and $c'=M-v_{K-1}$. All steps
except the last are unclipped: $D_{k+1}=(1+\theta)D_k+h$ for $k<K-1$, and
$D_K\le(1+\theta)D_{K-1}+h$. The node $z=M$ is in the $z$-grid, and
$F_M(v,M)=(M-v)^2$. Every correction is nonnegative and
$d_x(v)\ge w_x(v)^2/4$.

If $K=1$, the node $M$ has the adjacent interval $[c_x,M]$ of length at least
$M/2$, so $\beta\le Q(M,M)\le-M^2/16$, and
$M^2/16\ge\max\{\theta^2M^2/379,h^2/20\}$ since $\theta\le1/4$ and $h\le M$.

If $K\ge2$, then $Q(M,M)\le-c'^2/4$ and
$Q(v_{K-1},M)\le c'^2-a^2/4$. If $c'^2\ge a^2/5$ the first is at most
$-a^2/20$, otherwise the second is below $-a^2/20$; so $\beta\le-a^2/20$.
Now $a=h+\theta D_{K-2}\ge h$, and the recursion gives
$D_{K-2}\ge(D_K-(2+\theta)h)/(1+\theta)^2$. If $h\le M/16$, then with
$D_K\ge M/2$ and $\theta\le1/4$,
$D_{K-2}\ge(\frac M2-\frac{9M}{64})\cdot\frac{16}{25}=\frac{23M}{100}$, so
$a\ge\frac{23}{100}\theta M$. If $h>M/16$, then $a\ge h>M/16\ge\theta M/4$.
Hence $a\ge\max\{h,\frac{23}{100}\theta M\}$ and
$\beta\le-\max\{h^2,(0.23\,\theta M)^2\}/20$, which
implies~\eqref{eq:twocenters} because $0.23^2/20>1/379$.

\emph{Consequences.} Filtering keeps both minimizers
(Proposition~\ref{prop:filter}), so every box of FG is $[0,M]^2$, and every
stage is of the form above with $U\ge\OPT=0$; this gives the gap bound. A
stage with $U-\beta\le\varepsilon$ has $\theta\le\sqrt{379\varepsilon}/M$
and $h\le\sqrt{20\varepsilon}$, so every step is at most
$h+\theta M<24\sqrt\varepsilon$, and the side of length at least $M/2$ in each
coordinate needs at least $M/(48\sqrt\varepsilon)$ nodes. The table of the
bag $\{x,z\}$ has at least $M^2/(2304\varepsilon)$ entries. CT stops only
after some stage certifies its target, and an acceptance with
$U-\beta<1$ needs a stage with gap below one.
\end{proof}

The instance is trivial for other methods; the point is that the failure
comes from the single center, not from conditioning. Replacing hulls by
unions of retained cells, on uniform meshes, removes it at the price of a
$\sqrt n$ factor per coordinate.

\paragraph{Algorithm UC$(\varepsilon)$ (uniform cells).}
For each coordinate keep a finite set $\mathcal C_i$ of \emph{retained cells},
closed intervals of positive length; initially $\mathcal C_i=\{[\ell_i,u_i]\}$.
Let $J$ be the least $j\ge0$ with $\frac12nLs^24^{-j}\le\varepsilon$, start
with $\hat x=\ell$, $U=F(\ell)$, and for $j=0,1,\dots,J$ with $h_j=s2^{-j}$:
\begin{enumerate}[label=(\roman*)]
\item \emph{Refine.} Let $\Lambda_{ij}=\{\ell_i+kh_j:k\in\Z_{\ge0}\}$ for
$i\in I_C$ and $\Lambda_{ij}=\{\ell_i+\lceil kh_j\rceil:k\in\Z_{\ge0}\}$ for
$i\in I_Z$. A retained cell $[a,b']$ is split at the points of
$\Lambda_{ij}\cap(a,b')$; the resulting \emph{stage cells} have consecutive
nodes as endpoints, and $G_i$ is the set of all their endpoints. Let
$w_i(v)$ be the largest length of a stage cell with endpoint $v$, ignoring
integer cells of length one, and define $d_i$, $D$, $Q$, $\beta_j$ and $m_i$
as in~\eqref{eq:correction}--\eqref{eq:minmarginal} on $G=\prod_iG_i$.
\item \emph{Solve.} Compute $\beta_j$, a minimizer $y_j$ and all $m_i$ by
Lemma~\ref{lem:dp}; if $F(y_j)<U$, set $\hat x=y_j$, $U=F(y_j)$.
\item \emph{Filter.} Let $\mathcal C_i$ be the set of stage cells $[a,b']$
with $\min\{m_i(a),m_i(b')\}\le U$. Do not take hulls.
\end{enumerate}
Return $\hat x$, $U$ and $\beta_J$. The current domain at stage $j$ is
$X_j=X\cap\prod_i\bigcup\mathcal C_i$.

Since $\Lambda_{i,j-1}\subseteq\Lambda_{ij}$ (for integers,
$\lceil kh_{j-1}\rceil=\lceil 2kh_j\rceil$), every stage cell is either a
segment between consecutive points of $\Lambda_{ij}\cap[\ell_i,u_i]$ or such
a segment clipped at $u_i$. Hence a continuous stage cell has length at most
$h_j$; an integer stage cell has length at most $\lceil h_j\rceil$, and if
that length is at least two then $h_j>1$ and $\lceil h_j\rceil<2h_j$.
Consecutive points of $\Lambda_{ij}$ are at least $h_j$ apart for $i\in I_C$
and at least $\max\{1,\lfloor h_j\rfloor\}\ge h_j/2$ apart for $i\in I_Z$.
Each stage cell contains at most one point of $\Lambda_{i,j+1}$ in its
interior.

\begin{lemma}[Cells: interpolation and safe filtering]\label{lem:cells}
Assume only Definition~\ref{def:curvature}. At every stage, for every
$x\in X_j$ there is a random $Y\in G$ with $\E Q(Y)\le F(x)$ and
$Y_i\in\{a,b'\}$ whenever $x_i$ lies in the stage cell $[a,b']$. Hence
$\beta_j\le\min_{X_j}F$, and $\min\{m_i(a),m_i(b')\}\le F(x)$ for every
$x\in X_j$ with $x_i\in[a,b']$. Every global minimizer lies in every $X_j$,
and $\beta_J\le\OPT\le U$.
\end{lemma}

\begin{proof}
For each $i$ choose a stage cell $[a_i,b_i]$ containing $x_i$. If $x_i$ is
an endpoint put $Y_i=x_i$; otherwise let $Y_i=b_i$ with probability
$(x_i-a_i)/(b_i-a_i)$ and $Y_i=a_i$ otherwise, independently. If $i\in I_Z$
and $b_i-a_i=1$, then $x_i$ is an endpoint. Let
$Z^{k}=(Y_1,\dots,Y_k,x_{k+1},\dots,x_n)\in\bar X$. Conditionally on
$Y_1,\dots,Y_{k-1}$, the function $\psi(t)=F(Z^{k-1}+(t-x_k)e_k)-\frac{L_k}2t^2$
is concave on $[a_k,b_k]$, so Jensen's inequality gives
$\E F(Z^k)\le F(Z^{k-1})+\frac{L_k}2\Var(Y_k)$. Summing over $k$,
$\E F(Y)\le F(x)+\sum_k\frac{L_k}2\Var(Y_k)$. If $Y_k$ is rounded, then
$\Var(Y_k)=(x_k-a_k)(b_k-x_k)\le(b_k-a_k)^2/4$ while both possible values have
$d_k\ge L_k(b_k-a_k)^2/8$; if not, $\Var(Y_k)=0\le d_k$. So
$\E D(Y)\ge\sum_k\frac{L_k}2\Var(Y_k)$ and $\E Q(Y)\le F(x)$. Since
$Q(Y)\ge\beta_j$ and, when $x_i\in[a,b']$ is the chosen cell, $Q(Y)\ge
\min\{m_i(a),m_i(b')\}$, the two bounds follow. A point $x\in X_j$ removed
at stage $j$ has some coordinate whose stage cells containing $x_i$ were all
discarded, so $F(x)>U_j\ge U$; in particular minimizers are never removed,
and by induction $X_J$ contains $S$ and $\beta_J\le\OPT$.
\end{proof}

\begin{theorem}[Uniform cells]\label{thm:cells}
Assume Definition~\ref{def:curvature}. For every $\varepsilon>0$,
UC$(\varepsilon)$ returns a feasible $\hat x$ and $\beta\le\OPT\le
F(\hat x)\le\beta+\varepsilon$ with a certificate (all stage tables and
filtering decisions) whose validity uses only Definition~\ref{def:curvature}.
Suppose moreover that $S$ is finite, $r=\max_i\abs{\pi_i(S)}$, and
$F(x)-\OPT\ge g_S\dist(x,S)^2$ on $X$; put $\kappa_S=\max\{1,L/g_S\}$. Then
every grid of every stage has at most
\begin{equation}\label{eq:cellcount}
K_S=12\,r\,\bigl(2\sqrt{n\kappa_S}+1\bigr)
\end{equation}
nodes per coordinate. For an explicit rational quadratic, UC$(2^{-q})$ uses
$O\bigl(p(N+\abs{\mathcal A})K_S^p\bigr)\poly(I+q)$ bit operations, and
EX with UC in place of CT returns an exact minimizer in
$O\bigl(p(N+\abs{\mathcal A})K_S^p\bigr)\poly(I+\log\kappa_S)$ bit
operations. None of $g_S$, $S$, $r$ is an input.
\end{theorem}

\begin{proof}
\emph{Gap.} Every stage cell that receives a correction has length at most
$h_j$ (continuous) or below $2h_j$ (integer), so
$D(y)\le\sum_iL_i(2h_j)^2/8\le\frac12nLh_j^2$ for every $y\in G$. Hence
$U-\beta_j\le F(y_j)-\beta_j=D(y_j)\le\frac12nLh_j^2$, which is at most
$\varepsilon$ at $j=J$. Validity is Lemma~\ref{lem:cells}.

\emph{Count.} Fix a stage $j$ and a cell $[a,b']$ retained by step~(iii). One
endpoint $v$ has $m_i(v)\le U$; let $z\in G$ attain $m_i(v)$, so $z_i=v$.
By Lemma~\ref{lem:cells}, $\beta_j\le\OPT$, so
$U\le F(y_j)\le\OPT+\frac12nLh_j^2$ and
\[
F(z)-\OPT\le U+D(z)-\OPT\le nLh_j^2 .
\]
Set growth gives $\dist(z,S)^2\le nLh_j^2/g_S\le n\kappa_Sh_j^2$, so
$\abs{v-\sigma}\le\rho:=\sqrt{n\kappa_S}\,h_j$ for some $\sigma\in\pi_i(S)$.
The endpoints of stage cells are points of $\Lambda_{ij}$ or $u_i$. An
interval of length $2\rho$ contains at most $2\rho/\delta+1$ points of
$\Lambda_{ij}$, where $\delta=h_j$ (continuous) or $h_j/2$ (integer), plus
possibly $u_i$; each endpoint belongs to at most two stage cells. So at most
$2r(4\sqrt{n\kappa_S}+2)$ cells are retained, and since each contains at most
one new interior node, the stage-$(j+1)$ grid has at most
$3\cdot2r(4\sqrt{n\kappa_S}+2)=K_S$ nodes. Stage $0$ has two nodes per
coordinate.

\emph{Work and bit length.} By Lemma~\ref{lem:dp} each stage takes
$O(p(N+\abs{\mathcal A})K_S^p)$ table operations, and
$J=O(\log(nLs^2/\varepsilon))=O(I+q)$. Continuous nodes have the form
$\ell_i+ks2^{-j}$ and lie in $[\ell_i,u_i]$, integer nodes are integers, so
all nodes have denominators dividing $D2^j$ and bit length $O(I+j)$; factor
values, corrections and messages have bit length polynomial in $I+j$. For
exact output, Theorem~\ref{thm:transfer}(a) is reached at
$q\le\poly(I)+\log_2\kappa_S$, as in the proof of Theorem~\ref{thm:exact}.
\end{proof}

\begin{remark}[What grading buys]\label{rem:grading}
With $r=1$, Theorem~\ref{thm:cells} gives accuracy-independent grids of
size $O(\sqrt{n\kappa})$ without trial gradings or caps. The graded
single-center grids of Section~\ref{sec:growth} reduce $\sqrt n$ to
$\log n$, which is what makes Theorem~\ref{thm:approx} fixed-parameter
tractable in $(p,\kappa)$; Proposition~\ref{prop:twocenters} shows that
they pay for this with a dependence on a unique minimizer. Whether exact
output admits an $f(p,\kappa_S,r)\poly(I)$ bound when $S$ is finite, or any
accuracy-independent grid bound when $S$ is a continuum, is open.
\end{remark}

\begin{remark}[Relation to rational-witness results]\label{rem:np}
Corollary~\ref{cor:height} is the box case, with an explicit constant, of a
classical fact: a quadratic program with rational data that attains its
minimum has a global minimizer of polynomial encoding length, obtained from
a rational linear system \cite{Vavasis1990}; Vavasis used it to show that the
decision version of quadratic programming is in NP, and Del Pia, Dey and
Molinaro extended NP membership to mixed-integer quadratic programs with
unbounded integer variables \cite{DelPiaDeyMolinaro2017}. On a bounded
mixed box the integer coordinates have polynomial length, so that extension
is not needed. These results certify $\OPT\le t$. Deciding $\OPT\le t$ for
box QP is NP-complete, and NP-hard already at treewidth two
\cite{DelPiaKhajavirad2026}, so polynomial-size, polynomial-time checkable
certificates of lower bounds $\OPT\ge t$ for all instances would place an
NP-complete problem in coNP. The
certificates of Theorem~\ref{thm:exact} have size $f(p,\kappa)\poly(I)$:
they exist for bounded $(p,\kappa)$ and are checked without trusting
growth. The explicit constants $R$ and $V$ are what make the stopping test
computable.
\end{remark}

\subsection{Explicit polynomial factors}\label{sec:polynomial}

Now let each factor be an explicit list of rational monomials of degree at
most $d\ge2$, with $d$ fixed (or encoded in unary).

\begin{lemma}[Interval curvature certificate]\label{lem:intcurv}
For each $i$ write $\partial_{ii}F=\sum_m\gamma_mx^{\alpha_m}$ and let
$[\underline\mu_m,\overline\mu_m]$ be the range of the monomial $x^{\alpha_m}$
on $\bar X$. Put $\hat L_i=\sum_m\max\{\gamma_m\overline\mu_m,
\gamma_m\underline\mu_m\}$ and $L_i=\max\{\hat L_i,0\}$. Then the $L_i$ are
upper coordinate curvatures, and they are computed in time polynomial in $I$.
\end{lemma}

\begin{proof}
The monomial $x^{\alpha}$ is a product of functions $x_k^{\alpha_k}$ of
distinct coordinates, so its range on the box $\bar X$ is the product of the
intervals $\{t^{\alpha_k}:t\in[\ell_k,u_k]\}$, each of which is computed
exactly from $\ell_k^{\alpha_k}$, $u_k^{\alpha_k}$ and, for even $\alpha_k$
with $\ell_k<0<u_k$, the value $0$. Hence
$\gamma_mx^{\alpha_m}\le\max\{\gamma_m\overline\mu_m,\gamma_m\underline\mu_m\}$
on $\bar X$, and $\partial_{ii}F\le\hat L_i\le L_i$. For every $x\in\bar X$,
$t\mapsto F(x+te_i)-\frac{L_i}2t^2$ has second derivative
$\partial_{ii}F-L_i\le0$ on its interval, so it is concave. The endpoints are
powers of degree at most $d$ of input rationals, so their bit length is
$O(dI)$, and there are at most as many terms as input monomials.
\end{proof}

The bound can overestimate when monomials cancel; the condition number in
the results below is computed with the accepted $L$. Any sharper bound with a
certificate checkable in polynomial time may be used instead.

\begin{proposition}[Lattice stopping rule]\label{prop:lattice}
Suppose every coordinate is integer, and let $Q_0$ be a positive common
denominator of the coefficients of the objective (after substituting fixed
coordinates). If $\hat x\in X$, $\beta\le\OPT$ and
$F(\hat x)-\beta<1/Q_0$, then $\hat x\in S$. It suffices that
$F(\hat x)-\beta\le2^{-q}$ with $q=1+\lceil\log_2Q_0\rceil=O(I)$.
\end{proposition}

\begin{proof}
At integer points every monomial value lies in $Q_0^{-1}\Z$, so
$F(X)\subseteq Q_0^{-1}\Z$. If $F(\hat x)>\OPT$ then
$F(\hat x)-\OPT\ge1/Q_0>F(\hat x)-\beta$, contradicting $\beta\le\OPT$.
\end{proof}

\begin{corollary}[Polynomial factors]\label{cor:poly}
For fixed $d$, under quadratic growth~\eqref{eq:growth} with the accepted
curvature bounds, CT$(2^{-q})$ returns a certified $2^{-q}$-approximation in
$f_d(p,\kappa)(I+q+1)^{C}$ bit operations with $C$ absolute. If all
coordinates are integer, CT$(2^{-(1+\lceil\log_2Q_0\rceil)})$ returns an exact
minimizer, certified by Proposition~\ref{prop:lattice}, in
$f_d(p,\kappa)(I+1)^C$ bit operations.
\end{corollary}

\begin{proof}
The interpolation, filtering, contraction, localization and grid-size
arguments of Sections~\ref{sec:grids} and~\ref{sec:growth} use only
Definition~\ref{def:curvature} and growth, not the quadratic form. In the bit
analysis of Theorem~\ref{thm:approx} only factor evaluation changes: a
monomial of degree at most $d$ at nodes with denominators dividing $A_j$ has
denominator dividing $D_FA_j^d$, where $D_F$ clears the coefficients, and
if the nodes have bit length at most $\lambda$ its value has bit length
$O(I+d\lambda)$; corrections and
messages then have a common denominator dividing $8D_FA_j^{d}$ (with $D_F$
also clearing $L$). Evaluation of an explicit factor costs a polynomial in
its length and these bit lengths, for fixed $d$. The second statement
follows from Proposition~\ref{prop:lattice} with $q=O(I)$.
\end{proof}

\begin{example}[Limits of the polynomial extension]\label{ex:polylimits}
(a) $F(x)=x^3-6x$ on $[1,2]$ has $\partial_{xx}F=6x\le12=\hat L$ and
$F(x)-\OPT=(x-\sqrt2)^2(x+2\sqrt2)\ge3(x-\sqrt2)^2$, so it has growth $g=3$
and $\kappa=4$; but its minimizer $\sqrt2$ and value $-4\sqrt2$ are
irrational, so no rational exact output exists for continuous polynomial
coordinates. (b) $F(x)=x^4$ on $[-1,1]$ has a unique minimizer and no
$g>0$, so Lemma~\ref{lem:unique-growth} fails for polynomials.
(c) For $k\ge1$, $F(x,y)=x^2-\frac12x^{2^k}+y^2$ on $[0,1]^2$ has
$\partial_{xx}F\le2$, $\partial_{yy}F=2$ and growth $g=\frac12$, since
$x^{2^k}\le x^2$; its input has $O(k)$ bits if exponents are written in
binary, but $F(\frac12,0)$ has denominator $2^{2^k+1}$. Binary-encoded
exponents therefore need a separate evaluation analysis.
\end{example}
